Congressional apportionment is a step in the redistricting process
that must be performed exactly.
An error in apportionment --- giving a state one more or fewer seat
than it is legally entitled to --- propagates through the entire
redistricting pipeline: the number of districts drawn, their target
population, and the resulting map are all wrong.
For a redistricting platform to be usable in litigation, commission
proceedings, or policy analysis, its apportionment step must be
demonstrably correct.

The official U.S.\ congressional apportionment is computed by the
Census Bureau using the Huntington-Hill method mandated by
2~U.S.C.\ \S2a \citep{usc2a}.
The Census Bureau publishes its official seat counts within 9 months
of each decennial census.
Any redistricting platform claiming to implement U.S.\ apportionment
law must reproduce these official seat counts exactly.

\textsc{bisect-apportion} currently approaches this requirement through four
design choices, with Census-reference verification still identified as future
hardening:
\begin{enumerate}
  \item \textbf{Double-precision arithmetic}: Priority values are compared
    as \texttt{f64} floating-point, which provides numerically stable
    comparisons for the populations used in apportionment.
  \item \textbf{Output reproducibility}: Current tests verify deterministic
    output on synthetic fixtures and method properties.  Full Census Bureau
    string/hash regression remains a future hardening step.
  \item \textbf{Comprehensive test suite}: 99 tests across the crate provide
    coverage of unit formulas, synthetic examples, prime-factor compositor
    behavior, nesting helpers, spectral helpers, and split strategies; 24 tests
    specifically target apportionment methods and paradox detection.
  \item \textbf{Separate functions per method}: \texttt{huntington\_hill()}
    (for the official U.S.\ method) and \texttt{apportionment\_divisor()}
    (for Webster, Adams, and Jefferson via \texttt{RoundingRule}) enable
    direct comparison with minimal shared-code risk.
\end{enumerate}

This paper documents the implementation in sufficient detail for
replication, audit, and extension.
Sections~\ref{sec:architecture}--\ref{sec:other-methods} describe the
code architecture and method implementations.
Section~\ref{sec:verification} describes the verification approach.
Section~\ref{sec:tests} documents the test suite.
Section~\ref{sec:comparison} presents comparison-table targets and the current
evidence boundary.

\subsection{Why Four Implemented Divisor Methods?}

The crate implements the four divisor methods because practitioners and
researchers need counterfactual analysis:
\emph{what would the seat allocation be if Congress had adopted
Webster instead of Huntington-Hill?}
\emph{What if Adams were used?}
These questions arise in apportionment reform advocacy, in academic
analysis of congressional representation, and in state-level
apportionment decisions where different methods are legally permitted.
A unified divisor-method implementation enables these comparisons without
requiring practitioners to seek separate tools for each method.  Hamilton remains
important for J.5 paradox analysis, but it is not yet a public
\textsc{bisect-apportion} apportionment API.
